\chapter{Languages that make interpretation inspectable}
\label{ch:languages}

By this point the reader has seen why a rule can be incomplete before it reaches
code. The next question is how to write the rule so that a compliance reviewer,
an engineer and a verification tool can inspect the same meaning. A gift
restriction needs an exception. A consent decision needs a history. A contract
needs parties, resources and dates. One notation cannot make all of those
questions equally clear.

The languages in this chapter are therefore introduced by the problem each one
solves. The names matter less than the distinction they expose. We will return
to the SPI and marketing examples, write small pieces by hand, and say exactly
which part of the interpretation each language can carry and which part it
leaves to review.


% BEGIN SOURCE UNIT M04:00
\section{Choose the question before choosing the language}
\label{merge:M04:00}

\label{sec:languages-question}
The gift example needs a condition over classified facts. A consent example
needs a history of grants and withdrawals. A payment agreement may need parties,
assets and deadlines. A dispute about the meaning of an exception needs
alternative readings and their supporting reasons. These are related questions,
but a representation designed for one does not automatically solve the others.
The purpose of reviewing formal legal languages is to identify which structures
help make each question precise.

A formal legal language is a language in which statements about a legal or
normative domain have defined syntax and computational meaning. The word
``legal'' identifies the intended subject matter and useful abstractions. It
does not mean the language has legal authority, that a program written in it
is enforceable, or that translation into it establishes a correct reading.
Authority still comes from the relevant source and institutional decision.

% BEGIN REVISION ADDITION teach-M04-00
\begin{figure}[H]
\centering
\TeachingBranch{Which question is being answered?}{Decision with exceptions}{State, resources and duties}
\caption{Different legal questions call for different computational languages.}\label{fig:teach-M04-00}
\end{figure}

% END REVISION ADDITION teach-M04-00
For our product, a useful language must support at least three forms of
inspection. A reviewer needs to see how a condition corresponds to source
text. An engineer needs to know exactly what values and events the program
expects. A verification tool needs unambiguous semantics. A language that is
expressive enough to represent any program may make verification harder; a
language that permits only simple decision tables may fail to represent
important temporal or exception structure.

This chapter develops the relevant mechanisms through examples. It does not
choose a universal replacement for the bank's Java system. The recommended
architecture keeps a bounded executable representation, a richer interpretation
record and explicit adapters for specialized analyses. The relationship between
those representations must itself be checked. Adding more languages without
defining that relationship would create new translation boundaries rather than
redundancy.
% END SOURCE UNIT M04:00

% BEGIN SOURCE UNIT P-02-literature:00
\subsection{The research families behind the design}
\label{merge:P-02-literature:00}

\label{sec:literature}

Computational law has developed several answers to different parts of this
problem. Logic programming makes statutory dependencies executable. Languages
such as Catala organize definitions and exceptions close to legal text. Normative
languages represent duties, permissions and violations. Contract languages model
interactions through time. Recent language-model research attempts to cross the
remaining boundary from prose to these representations. A useful product combines
these contributions selectively, because their guarantees concern different
objects.
% END SOURCE UNIT P-02-literature:00

% BEGIN SOURCE UNIT P-02-literature:01
\subsection{Scope and method of this review}
\label{merge:P-02-literature:01}


The review began with the seven requested links and expanded through their
references, forward-citation searches for Catala and Stipula, and searches on
executable law, legal reasoning, formalization, dates, testing and process
compliance. ResearchAssistant supplied discovery records and local PDF extraction.
The saved OpenAlex citation searches returned 52 entries for the selected Catala
record and 13 for the selected Stipula journal record. Those are retrieval counts
for particular indexed editions, not complete or comparable measures of influence.
Preprints, revised editions and new publications can be missing or duplicated.

The initial proposal retained 38 paper editions representing 37 distinct
works. The first unified library retained 50 editions representing 49 works,
including the later legal-search and uncertainty studies. The present
audit adds two legal-reliability studies, bringing the academic collection
to 52 editions representing 51 works. Regulatory documents, standards
and software records are retained alongside them.
The AAAI tax-reasoning paper and its extended preprint are retained separately;
all seven requested sources are present. Technical reading covered the relevant
methods, semantics, evaluation and supporting appendix material. The accompanying
notes identify the inspected sections, version differences, implementation checks
and limits of adoption \citep{libraryreview}. This is a focused technical review
supporting a product decision, rather than a claim to have enumerated every
citation. Unavailable secondary leads are recorded separately and do not support
implementation claims.

% BEGIN REVISION ADDITION teach-P-02-literature-01
\begin{figure}[H]
\centering
\TeachingFlow{Primary technical method}{Assumptions and evaluation}{Scoped borrowing for LegalMath}
\caption{Literature review must inspect the mechanism being borrowed.}\label{fig:teach-P-02-literature-01}
\end{figure}

% END REVISION ADDITION teach-P-02-literature-01
The expanded search treats eight families separately: executable statutes;
defaults and defeasible priorities; obligations and violations; temporal process
compliance; authority and provenance; legal languages, compilers and testing;
neurosymbolic formalization and evaluation; and deployable rule engines. Important
foundations predate Catala, while contract languages have partially separate
citation communities. The retained forward-citation records have individual
dispositions in \path{docs/papers/coverage.md}; additional searches and backward
references supply the missing mechanisms. Selection favors primary sources that
define a needed mechanism or report an informative failure. Downloading a paper
or screening its abstract is distinguished from inspecting its technical method.
This is a bounded review of the prototype's design questions. Remaining retrieval
gaps are stated rather than counted as read papers.

Version control is particularly important here. The ARc paper was first posted
in 2025, but the inspected second version is from July 2026. The eFLINT paper's
inspected version is from August 2026. The Stipula preprint requested by the
sponsor has three authors; the later journal edition has a different author list.
Bibliographic records and local filenames make these distinctions explicit.
Published results are reported as their authors' findings; none of their
benchmarks or proof developments was independently rerun for this proposal.
% END SOURCE UNIT P-02-literature:01

% BEGIN SOURCE UNIT M04:01
\section{Rules do not all behave like ordinary implications}
\label{merge:M04:01}

\label{sec:languages-norms}
An ordinary material implication says that if a proposition is true, another
proposition is true. A legal norm often says that, under specified circumstances,
someone ought to do something. The existence of the duty does not establish
that the action occurred. A rule requiring a document to be supplied cannot
be used as evidence that the document was supplied. The distinction between
facts and normative outputs is central to input/output logics
\citep[secs.~2--4]{inputoutput2000}.

For a simple illustration, suppose an adopted specification says that a
distributor must provide a named warning when a condition holds. The condition
can create a duty record. Performance requires separate evidence of the
warning being delivered to the correct recipient in the required context.
If the program inserts the duty's content into its factual database as though
it had happened, every obligation can appear automatically satisfied.

Prohibitions and permissions also need care. Failure to derive a prohibition
may mean the act is permitted under one explicitly defined logic; in a bank
workflow it may instead mean the relevant sources or facts are missing. The
literature distinguishes strong and weak permissions under particular
defeasible-logical definitions \citep[secs.~2--4]{permissions2012}. Those
definitions are not a general license to treat an empty query result as
authorization. The selected RuleIR gift result remains a subcondition, with
permission supplied only by a separate approved decision policy.

% BEGIN REVISION ADDITION teach-M04-01
\begin{figure}[H]
\centering
\TeachingBranch{General rule applies}{Exception defeats its conclusion}{No exception established}
\caption{Defeasible reasoning makes the exception relationship explicit.}\label{fig:teach-M04-01}
\end{figure}

A general fee rule and a special discount exception cannot be combined by treating every sentence as an unconditional implication. The formalization must state when the exception applies and what happens when its applicability is unknown.


% END REVISION ADDITION teach-M04-01
Many legal rules are defeasible: they normally support a conclusion, but an
exception or stronger opposing rule can override that support. A defeasible
rule is not the same as an ordinary implication with an informal comment
attached. Its evaluation must account for applicable opposing rules and
priorities. Defeasible-logic systems formalize different ways to do this,
including rules that support conclusions and defeaters that block opposing
support \citep[sec.~2]{defeasible2000}.

The engineering implication is simple but demanding. ``Has exceptions'' is
not a complete semantics. The specification must say what happens when no
exception is known, one is established, two compete, or an exception remains
unknown. It must also say whether a priority is supplied by legal authority,
an explicit interpretation decision or a programming convention. File order
is not an adequate hidden source of legal priority.

The early logic-program treatment of the British Nationality Act demonstrated
how a selected body of legislation could be decomposed into rules with
inspectable reasoning \citep{sergot1986}. That work is valuable here as an
example of disciplined formalization, including attention to time and negative
information. It should not be read as a claim that ordinary legal prose can
always be converted automatically into a complete logic program.
% END SOURCE UNIT M04:01

% BEGIN SOURCE UNIT P-02a-foundations:00
\subsection{Why exceptions, duties and permissions need different logics}
\label{merge:P-02a-foundations:00}


A threshold calculation asks whether a proposition follows from evidence. A duty
asks what an actor ought to do in those circumstances. Treating both as ordinary
implication loses the possibility of noncompliance: an instruction to conduct an
annual review is not evidence that the review occurred. Input/output logic makes
this separation explicit. A normative system transforms factual conditions into
required outputs, without requiring every input to be an output or permitting
unrestricted contraposition \citep[secs.~2--3]{inputoutput2000}.

Let $A$ contain factual propositions and let $G$ contain pairs $(a,x)$ expressing
that condition $a$ calls for outcome $x$. Writing $\operatorname{Cn}$ for classical
consequence, the paper's simplest operation is
\begin{equation}
 \operatorname{out}_1(G,A)
 =\operatorname{Cn}\!\left(G\!\left(\operatorname{Cn}(A)\right)\right).
 \label{eq:io}
\end{equation}
Here $G(S)=\{x\mid (a,x)\in G\text{ and }a\in S\}$ collects the required
outputs whose conditions are in $S$. Thus the two uses of $G$ denote the same
normative system, first as a set of pairs and then as the operation of applying
those pairs. First derive the input conditions, then apply the normative pairs, then unpack
their consequences. Section 3.2 characterizes this operation by strengthening
inputs, conjoining outputs and weakening outputs. A pair connecting ``review
due'' with ``review performed'' supplies a required outcome; it does not create
a recorded fact of performance. Other variants change disjunction and reuse.
LegalMath borrows this separation rather than implementing all input/output logics.

% BEGIN REVISION ADDITION teach-P-02a-foundations-00
\begin{figure}[H]
\centering
\TeachingCompare{Permission}{Obligation}{Reparation after violation}
\caption{Permission, duty and the consequence of breach are different normative relations.}\label{fig:teach-P-02a-foundations-00}
\end{figure}

% END REVISION ADDITION teach-P-02a-foundations-00
Exceptions ask a different question: which applicable rule defeats another?
\citet[sec.~2]{defeasible2000} distinguish facts, strict rules, defeasible rules,
defeaters and an explicit superiority relation. Their proof theory records both
definite and defeasible provability, and explicit failure to establish either.
A defeasible conclusion needs supporting premises, failure of definite contrary
proof and a response to every applicable opposing rule. Different supporting
rules may defeat different opponents. This ``team defeat'' differs materially
from selecting one table row or finding one Catala exception.

An unretrieved agreement is therefore not a defeater merely because the database
has no row. The product must distinguish explicit negative evidence, a
closed-world assumption and a reviewed default. LegalMath 0.1 chooses nested,
source-linked defaults with conflicts between competing sibling exceptions. It
does not implement the full defeasible proof theory. General priorities remain
an extension requiring a stated defeat policy and conformance tests.

Permissions add another distinction. In this proof theory, a weak permission
requires a derivation establishing that the opposite obligation is not
derivable; an unfinished search or an empty database response is insufficient.
A strong permission has explicit normative support. \citet[secs.~2--3, 6]{permissions2012} distinguish permissive
rules from defeaters and distinguish obligation from factual fulfillment.
Their algorithms transform a finite theory while preserving its defeasible
extension. The soundness, completeness and linear-complexity results in section 5
concern that theory and specified data-structure assumptions, not arbitrary
regulations. For the bank, failure to derive a prohibition must never by itself
become a positive trade authorization.

A remedy is different again. The logic of violations of
\citet[secs.~2--5]{violations2006} uses an ordered reparation connective: perform
the primary duty; if it is violated, a further duty can become relevant.
Replacing that structure with ordinary disjunction can make the remedy appear
an equally acceptable way to satisfy the original requirement. A warning delivered
after a deadline can be useful remediation without making the original control
timely. LegalMath retains the breach assessment and records subsequent performance
separately; general compensation and sanction reasoning remain outside its first
event profile.
% END SOURCE UNIT P-02a-foundations:00

% BEGIN SOURCE UNIT M04:02
\section{Controlled English: readable structure with restricted freedom}
\label{merge:M04:02}

\label{sec:languages-logicalenglish}
Ordinary English permits pronouns, elliptical references, implied subjects and
several possible attachments of a qualification. Those features make writing
compact, but they leave a translator with choices. Controlled English restricts
some of that freedom so a sentence can be mapped into logic by explicit rules.
Logical English develops this approach for law and education
\citep[language specification and examples]{logicalenglish2023}.

Consider the sentence ``If it is a gift and is linked to the product or its
type, it is prohibited unless discounted.'' The sentence is short, but ``it''
could refer to the offer or component, ``its'' needs a binding, and ``discounted''
does not identify the exception. A controlled rendering should repeat the
subject where necessary and name each definition. Repetition in the language
is useful if it removes a legally consequential ambiguity.

One possible authoring form for our own restricted profile is:
\begin{lstlisting}[caption={Illustrative LegalMath authoring notation. This is a proposed template, not claimed to be valid Logical English or Catala source.}]
subject: component C in offer O
scope:
  actor(O) satisfies distributor_profile
  activity(O) satisfies authorised_fund_promotion_profile
condition:
  C satisfies gift_definition
  AND (
    C has specific_product_link
    OR C has product_type_link
  )
  AND NOT (C satisfies fee_discount_definition)
consequence:
  selected_gift_restriction_triggered
\end{lstlisting}
The form makes four decisions inspectable: the subject, the scope, the
condition and the named consequence. Every definition identifier refers to
a versioned record. A reviewer can ask why a definition applies without
reverse-engineering the syntax tree. An engineer can translate the connectives
deterministically once the template is fixed.

% BEGIN REVISION ADDITION teach-M04-02
\begin{figure}[H]
\centering
\TeachingFlow{Restricted English syntax}{Explicit variables and rules}{Executable interpretation}
\caption{Controlled English reduces freedom to improve inspectability.}\label{fig:teach-M04-02}
\end{figure}

% END REVISION ADDITION teach-M04-02
Variable binding is a practical source of error. If one occurrence of $C$
refers to the voucher and another to the whole offer, the sentence has changed
its subject midway. The authoring language should enforce a declared variable
and type environment. A reference to an undeclared subject should be rejected.
An association such as ``component C belongs to offer O'' should be explicit
when the rule uses facts about both.

Negation requires similar precision. ``There is no evidence that C is a gift''
is not the same assertion as ``evidence establishes that C is not a gift.''
Logical-programming systems may distinguish explicit negation from failure to
prove a proposition. Our evidence model records known false, unknown and
conflict separately. A controlled sentence that asks about missing evidence
should use that information state explicitly; it should not rely on a model
to guess which sense of negation the author intended.

Controlled English improves the boundary between reviewed meaning and
implementation. It does not remove the earlier task of choosing the controlled
sentence from the original circular. A reviewer can still choose the wrong
definition or omit a phrase. That is why the product retains source anchors,
alternative templates and contrasting scenarios alongside the readable rule.
The language makes a chosen interpretation precise enough to challenge.
% END SOURCE UNIT M04:02

% BEGIN SOURCE UNIT P-02-literature:02
\subsection{Statutory programs, literate rules and collaboration evidence}
\label{merge:P-02-literature:02}


The British Nationality Act logic program is an early demonstration that selected
legal provisions can be represented as executable dependencies. Its significance
for this project is architectural: a conclusion can be explained through the
conditions supporting it, while rules and facts are represented separately.
The work also confronts negative information and temporal conditions, which
remain central to bank data. It does not show that arbitrary legal prose has a
unique automatic translation \citep{sergot1986}.

The literature turns that mechanism into a review question. If a source packet
contains a base rule and two exceptions, the reviewer should see the dependency
from each exception to the rule, the source span supporting each guard, and the
result when both guards fire. Catala's default calculus makes this failure
visible: two applicable exceptions can conflict even when their computed values
agree \citep[secs.~3--4]{catala2021}. The point for LegalMath is the traceable
comparison, not another automatic priority rule.

For LegalMath, this suggests a representation in which an exception names the
rule it qualifies and preserves the source supporting its priority. A discount
exception to a restriction on gifts should remain visibly attached to that
restriction. Two independently imported exceptions that both apply should trigger
a review of their relationship, not an arbitrary choice by the language model.
The product should also distinguish this normative priority from a document's
publication date: a newer explanatory document does not mechanically override
every older source it mentions.

% BEGIN REVISION ADDITION teach-P-02-literature-02
\begin{figure}[H]
\centering
\TeachingCompare{Statutory logic program}{Source-linked rule authoring}{Observed collaboration practice}
\caption{A formal language and evidence that people can use it answer different questions.}\label{fig:teach-P-02-literature-02}
\end{figure}

% END REVISION ADDITION teach-P-02-literature-02
The Catala compilation result is relevant but narrower than the phrase
``verified legal software'' might suggest. The paper proves properties of a
translation from the core default calculus into a lambda calculus with exceptions,
with a mechanized development in F*. It distinguishes that model from the
separately implemented production compiler. The inspected repository's
formalization notes additionally disclose an admitted supporting lemma; this
review has not replayed the development \citep[sec.~4]{catala2021,catalacode}.
The useful precedent is to define a small semantics and state exactly which
translation preserves it. It provides no theorem that a human or model has
captured the meaning of an SFC circular.

Catala's literate presentation also addresses the organizational problem.
\citet[sec.~4.2]{huttner2022} describe collaboration between legal experts and
programmers around a common representation. The legal expert can challenge a
condition without reverse-engineering application code, and the engineer can
make an ambiguity concrete. The exploratory lawyer study in the original Catala
paper gives useful interface evidence, but its small educational setting does
not establish reliable defect detection by bank reviewers
\citep[sec.~6]{catala2021}. LegalMath should borrow the collaboration model and
test its own interface with the intended users.

Controlled language offers another way to make formal statements reviewable.
Logical English retains a readable rule syntax while assigning a computational
meaning to predicates and variables \citep{logicalenglish2023}. For this product,
a controlled paraphrase should be generated from the adopted rule, so changing
the rule changes its explanation. It should display how variables are bound:
``the client'' must refer to the same person throughout a condition. Readability
cannot substitute for a defined semantics, and unconstrained English should not
be presented as if it were executable merely because it looks regular.
% END SOURCE UNIT P-02-literature:02

% BEGIN SOURCE UNIT M04:03
\section{Catala: defaults, exceptions and source-linked authoring}
\label{merge:M04:03}

\label{sec:languages-catala}
A source-linked exception needs two things at once: a defined value when a
guard applies and a record of which text gives that guard authority. Catala
joins those concerns by placing its exception calculus beside the legal text
\citep[secs.~2--4]{catala2021}. The question for this section is how a bank
profile should expose competing guards; the example below treats their nesting
and evidence as observable cases.

The useful structural question is whether an exception is nested inside a
qualification or sits beside another exception. A nested exception is considered
only after its parent condition is active; two sibling exceptions are competing
answers to the same base rule. Moving a condition one level can therefore change
which evidence is needed before a value is available, even when every word of the
policy remains present.

Take a base value $b=\mathsf{false}$ and two sibling guards $g$ and $h$, each of
which returns $\mathsf{true}$. If $g$ is established and $h$ is refuted, the
exception supplies the value. If $g$ is established and $h$ is unknown, the
result remains unresolved because the second exception may still apply. If both
are established, the result is a conflict even when their values agree. Catala's
empty-applicable-value and conflict distinction motivates this treatment; the
bank profile adds the separate information state for incomplete evidence. Equal
Booleans therefore do not settle the priority question.

% BEGIN REVISION ADDITION teach-M04-03
\begin{figure}[H]
\centering
\TeachingBranch{Default value}{Applicable exception changes it}{Competing exceptions need treatment}
\caption{Catala makes default and exception structure part of the program.}\label{fig:teach-M04-03}
\end{figure}

% END REVISION ADDITION teach-M04-03
LegalMath borrows explicit exception structure but adds its own treatment of
incomplete evidence. In the present RuleIR profile, every sibling guard is
evaluated. Two true guards conflict. One true guard with an unknown competitor
remains unknown because the competitor may become active. One true guard with
all competitors false selects its value. When all guards are false, the base
is selected. This is a separately specified design; it is not a claim that
RuleIR implements the whole Catala calculus.

The distinction matters for ensemble comparisons. A Catala model and a RuleIR
model can be useful independent representations, but they do not become
equivalent merely because both have an ``exception'' construct. The adapter
must define how missing evidence, no applicable value and competing exceptions
are represented. A comparison outside that common fragment should return
unsupported or expose the difference, not choose a convenient interpretation.

Catala's compiler work also illustrates how formal claims should be scoped.
The paper establishes properties for a core translation; the production
compiler and surrounding tooling are separate implementation objects. The
retained repository inspection also records an admitted supporting lemma in
the formalization notes. This project has not replayed those proofs. We borrow
the explicit language structure and preservation discipline, while requiring
independent evidence for any concrete adapter used in the bank product.

\citet[sec.~4.2]{huttner2022} discuss legal/programmer collaboration as a way
to keep source meaning and executable rules connected. The useful lesson is
organizational as well as technical. A legal reviewer can challenge an
exception's interpretation while an engineer checks that its computational
consequence matches that interpretation. Neither role should be expected to
certify the other's entire domain from a single code review.
% END SOURCE UNIT M04:03

% BEGIN SOURCE UNIT P-01a-language-lesson:02
\subsection{Applying the exception model to SPI explanations}
\label{merge:P-01a-language-lesson:02}


Apply that distinction to the SPI explanation requirement. In Annex 1
paragraph 10.3, the adopted profile makes the explanation guard active when the
client asks for one or raises material queries. Catala's source-linked style
motivates keeping that guard beside the text \citep[secs.~2--4]{catala2021},
while the banking profile must still specify how missing evidence is represented.

Write the guard explicitly rather than splitting the two conditions into
competing exceptions:
\begin{lstlisting}[caption={An exception with an explicit guard. The combined guard avoids creating two competing exceptions for the same requirement.}]
normally: product_explanation_required = false
exception when client_requests_explanation
               OR client_has_material_queries:
    product_explanation_required = true
always retain: current_offering_documents_required = true
\end{lstlisting}
The \emph{guard} is the condition after ``when''. If the guard is true, use the
exception; if false, use the general value. If evidence about the guard is
missing, LegalMath leaves the requirement unresolved. The separate offering-
document requirement is not inside the explanation exception. A programmer
who implements one switch for both has changed the model
\citep[Annex 1, para.~10.3]{sfcspiannex1}.

% BEGIN REVISION ADDITION teach-P-01a-language-lesson-02
\begin{figure}[H]
\centering
\TeachingBranch{General financial assessment}{Named exception or qualification}{Explanation of the selected result}
\caption{An explanation should show why the default or exception was selected.}\label{fig:teach-P-01a-language-lesson-02}
\end{figure}

% END REVISION ADDITION teach-P-01a-language-lesson-02
The generic distinction now has a concrete test. In the SPI packet, the request
flag is true and the material-query flag is false, so the combined guard selects
the exception. If the request flag is true while the query flag is unknown, the
result is unresolved: a separate sibling exception may still be active. If both
flags are true, two separately written exceptions conflict even when both return
``explanation required''. A single guard containing their OR expresses one
reviewed exception and avoids that sibling conflict.

This is the difference between borrowing Catala's organization and implementing
Catala itself. Catala distinguishes an empty applicable value from a conflict;
LegalMath adds incomplete bank evidence as an explicit unknown state. Its
proposed default operator evaluates every sibling guard, uses nesting for an
adopted priority, and never treats file order as legal priority. The earlier
SPI-Demo1 Java example uses the equivalent explicit Boolean condition and rejects
the general default operator. The subsequently completed v0.1 runtime implements
that operator with its own conformance tests; the two compiler profiles remain
distinct.
% END SOURCE UNIT P-01a-language-lesson:02

% BEGIN SOURCE UNIT M04:04
\section{Metadata is part of the meaning: LegalRuleML}
\label{merge:M04:04}

\label{sec:languages-legalruleml}
A formula alone does not say which provision it interprets, which jurisdiction
it concerns, when it applies, or whether a competing formula is also
defensible. LegalRuleML addresses representation of legal rules together with
such contextual and provenance information. The work on legal interpretations
explicitly represents alternative readings associated with source material
\citep[secs.~2--4]{legalinterpretations2014}; the later OASIS specification
provides a broader standardized vocabulary \citep{legalruleml2021}.

For the fee-discount issue, two candidate rules can point to the same source
span while naming different interpretation records. One record might restrict
the exception to an initial reduction; another might include a documented
refund. Their source identity is shared, but their assumptions and consequences
differ. The system should not assign a single interpretation identifier to
both merely because the source paragraph is the same.

% BEGIN REVISION ADDITION teach-M04-04
\begin{figure}[H]
\centering
\TeachingFlow{Rule content}{Authority, jurisdiction and time}{Applicable contextualized rule}
\caption{LegalRuleML identifies context that a bare formula does not express.}\label{fig:teach-M04-04}
\end{figure}

% END REVISION ADDITION teach-M04-04
The representation must also preserve the authority of a selection. A legal
reviewer may adopt one interpretation for a defined product and date range.
That selection is an event with an author, reason and scope. It does not erase
the alternative or establish that the alternative was syntactically invalid.
A future source clarification may require reconsidering the choice. The
provenance record is what lets the system determine which decisions depend on it.

LegalRuleML is a representation standard, not a turnkey engine for resolving
every legal dispute. A serialized permission or priority still requires a
specified evaluation semantics. Different reasoning engines may support
different fragments or make different assumptions. For the prototype we should
borrow the discipline of explicit context and alternatives without requiring
the entire standard to become the first runtime dependency.

This leads to a useful separation in the proposed storage model. A source span
identifies words. An interpretation claim identifies an adopted or proposed
meaning. A formal rule identifies a computation. A review decision identifies
who accepted which claim under which context. Keeping those identities
separate prevents a hash of the source text from being used as if it also
identified the meaning chosen from that text.
% END SOURCE UNIT M04:04

% BEGIN SOURCE UNIT M04:05
\section{Stipula: a contract as parties, resources and transitions}
\label{merge:M04:05}

\label{sec:languages-stipula}
Stipula becomes relevant when a legal arrangement changes through actions over
time. Its original paper describes agreements, parties, fields, assets,
functions, states and scheduled events \citep[secs.~3--6]{stipula2021}. A
function can be available only to a specified party in a specified state, with
a precondition on the current data. Executing it changes fields or resources
and moves the contract to another state. This is a much richer object than
a single decision-table row.

To understand the mechanism, consider a synthetic service agreement. It begins
in a state awaiting agreement on terms. After the designated parties agree,
it enters an active state. An authorized party can request a permitted service
only while the agreement is active. A withdrawal changes the state so that
the same request is no longer available under that agreement. A timed event
may move the agreement into an expired state. This example illustrates a state
model; it is not a claim about the legal effect of any particular bank consent.

The state matters because the same action name can have different consequences
at different points in the history. A function called \texttt{requestService}
cannot be understood by its input parameters alone if its availability depends
on whether consent is active. A current Boolean copied from a stale database
may also be inadequate. The state should be derived from an attributable
history or a reviewed snapshot plus a complete subsequent event stream.

Stipula's agreement construct explicitly identifies who participates and who
must agree on the initial terms. This is useful as a modeling discipline:
the identity of the party agreeing is not an untyped annotation on a field.
For LegalMath, the analogous review question is who has authority to approve a
meaning, data mapping or release. We borrow the importance of explicit parties
and transitions, not a claim that a regulatory circular is itself a private
contract between the bank and regulator.

% BEGIN REVISION ADDITION teach-M04-05
\begin{figure}[H]
\centering
\TeachingFlow{Parties hold resources}{Permitted transition transfers them}{New contractual state}
\caption{Stipula models contracts through state and controlled resource movement.}\label{fig:teach-M04-05}
\end{figure}

% END REVISION ADDITION teach-M04-05
Assets in Stipula are distinguished from ordinary fields. An ordinary numeric
field can be copied as information; an asset transfer must account for the
resource leaving one location and arriving at another. In the original
semantics, asset movement prevents creating a negative quantity and preserves
the relevant resource through transfer. This supports questions about duplicate
spending, locked resources and the availability of contractual actions. It is
not a model of market liquidity or a substitute for the bank's general ledger.

For a synthetic escrow, suppose a contract holds a resource amount $b$ and
transfers $q$ to a recipient. A local arithmetic description is
\begin{equation}
 b'=b-q,\qquad r'=r+q,\qquad 0\leq q\leq b,
 \label{eq:asset-transfer}
\end{equation}
where $r$ is the recipient's modeled holding. Adding the first two equalities
gives $b'+r'=b+r$. This is a conservation property of the modeled transfer.
It says nothing about whether the underlying asset exists off-system or whether
the transfer is legally authorized. Those are separate premises. The example
shows how a formal language can support a useful exact property without
claiming to settle the whole legal arrangement.

Scheduled events add another dimension. A Stipula event can schedule a
continuation at a time and execute it only if the contract remains in a
specified state. An action taken before the event can change whether that
continuation is applicable. The ordering of an event and a function call at
the same time is consequently significant. The inspected original semantics
gives scheduled events precedence in the relevant equal-time situation. A bank
adapter must not assume that this is the legally correct ordering for its own
deadlines. The source, clock precision and event-order authority need review.

This example also distinguishes observation from permission. A formal contract
may describe which function calls its own execution platform permits. A bank
audit system must still record actions that occurred outside a permitted
workflow. Deleting an observed transaction because the policy would have
forbidden it would erase evidence of a possible violation. A normative monitor
therefore needs a clear separation between advice about an action and recording
its actual occurrence.
% END SOURCE UNIT M04:05

% BEGIN SOURCE UNIT P-01a-language-lesson:03
\subsection{Applying the state model to client consent}
\label{merge:P-01a-language-lesson:03}

\label{sec:stipula-tutorial}

A wealth calculation is a function of facts. An agreement also has a history.
Before signature, after signature, after withdrawal and after an amendment, the
same attempted action can have different consequences. Stipula is a language
for expressing those interactions. The requested 2021 paper defines parties,
agreement terms, fields, assets, states, guarded functions and scheduled events
\citep[secs.~3--4]{stipula2021}. The original language in that paper is
type-free; the typed decision representation proposed here is a separate design.

Read the rental example as a state machine rather than as a second glossary.
The party identifies who may act; the agreement binds the terms; fields hold
ordinary information; assets carry conservation obligations; the state gates
which actions are available; a guarded function changes fields or resources and
moves the state; and a scheduled event attempts a transition later if its state
condition still holds. This mapping matters because a function's name alone
cannot tell the reader whether a withdrawal, expiry or payment has changed what
that function means.

For example, the rental contract in the paper first records an offer. The
borrower accepts by supplying the required payment, which moves the contract
into its using state. The contract provides the use information and schedules
an event for the end of the rental period. Either an appropriate borrower action
or the scheduled event can finish the interaction and release the held payment.
The language gives a precise account of which party can trigger each step and
when the payment moves. Assets have special transfer behavior; they are not
ordinary numeric fields that an assignment may freely duplicate
\citep[sec.~3, rental example and grammar]{stipula2021}.

% BEGIN REVISION ADDITION teach-P-01a-language-lesson-03
\begin{figure}[H]
\centering
\TeachingFlow{Consent granted}{Consent withdrawn}{A later action sees the new state}
\caption{A consent history is a state problem, not a permanent client flag.}\label{fig:teach-P-01a-language-lesson-03}
\end{figure}

% END REVISION ADDITION teach-P-01a-language-lesson-03
The banking analogue uses the state and actor ideas but does not transfer client
assets. The following is explanatory pseudocode, deliberately distinguished
from compilable Stipula syntax:
\begin{lstlisting}[caption={A consent model inspired by Stipula. ``Record'' observes an event; ``allow'' describes the control's decision.}]
parties: Client, Bank
terms: client_id, product_category, threshold, monitoring_mode
states: NOT_ESTABLISHED, ACTIVE, WITHDRAWN

NOT_ESTABLISHED:
    record Client.written_agreement(terms, evidence)
    when Bank.confirmed_required_explanations(evidence)
    -> ACTIVE

ACTIVE:
    record Client.withdrawal(instruction_id) -> WITHDRAWN
    allow Bank.apply_streamlining(transaction)
    when category_matches(transaction, terms)
         AND all_other_required_conditions_hold(transaction)

WITHDRAWN:
    do not allow Bank.apply_streamlining(transaction)
    record a new, fully evidenced agreement -> ACTIVE
\end{lstlisting}
The pseudocode makes two boundaries testable. The terms bind consent to a
client, category and arrangement, so a new agreement after withdrawal creates a
new consent generation rather than deleting the withdrawal. A transaction
carrying the old generation must remain distinguishable from one carrying the
new one. The regulatory duties still come from the applicable sources: client
and bank agreement can establish the consent contemplated by those sources, but
it cannot amend an SFC requirement. The figure therefore supports a state-history
test, while the source decision remains a separate review object.

\begin{figure}[htbp]
\centering
\begin{tikzpicture}[node distance=17mm, every node/.style={font=\small},
box/.style={draw=teal,fill=pale,rounded corners,align=center,text width=33mm,minimum height=14mm},
arr/.style={-{Latex[length=2mm]},thick,draw=navy}]
\node[box] (none) {Consent not established};
\node[box,right=of none] (active) {Active consent\\for these terms};
\node[box,right=of active] (withdrawn) {Withdrawn};
\draw[arr] (none)--node[above=5pt,align=center,font=\footnotesize]{Evidenced\\agreement}(active);
\draw[arr] (active)--node[above=5pt,font=\footnotesize]{Withdrawal}(withdrawn);
\draw[arr] (withdrawn.south) to[bend left=35] node[below]{New evidenced agreement} (active.south);
\end{tikzpicture}
\caption{Consent transitions for the example. An incomplete event history produces
an unresolved assessment outside these established states. It is not evidence
of an active agreement.}
\label{fig:consent-tutorial}
\end{figure}

Execution requires a state, stored terms and a history of ordered events. In
Stipula's operational semantics, the runtime also maintains pending scheduled
events and a clock. In the inspected 2021 semantics, due events take priority
over function calls at the same time. That is a language rule, not a rule of
Hong Kong banking law. LegalMath must adopt an explicit ordering policy for
withdrawal and order events, record the authority for their sequence, and retain
an unresolved result where the ordering cannot be established. It cannot import
Stipula's clock convention as if the circular prescribed it.

There is a second adaptation. A control may prohibit streamlined treatment after
withdrawal, but the bank might nevertheless have processed a transaction.
Deleting that transaction from the history because it is an illegal transition
would hide the very incident a compliance system needs to detect. LegalMath
therefore separates the observed event record from the decision about whether
the proposed action is permitted by the selected control. Normative languages
such as eFLINT and Symboleo help make this distinction explicit: a duty or
violation is represented separately from the fact that an action happened
\citep{eflint2020,symboleo2020}.
% END SOURCE UNIT P-01a-language-lesson:03

% BEGIN SOURCE UNIT M04:06
\section{A small state contract that a reader can execute by hand}
\label{merge:M04:06}


A complete miniature model helps distinguish the legal-language idea from its
implementation in a particular tool. Consider a synthetic escrow agreement in
which a payer deposits a fixed amount, a recipient may receive it after an
accepted delivery, and the payer can recover it after an expiry if delivery has
not been accepted. The example is chosen to teach states, parties, resources and
time. It is not a model of an actual bank product or a proposed interpretation of
the gift circular.

Let the contract state be one of \texttt{UNFUNDED}, \texttt{HELD},
\texttt{PAID} and \texttt{REFUNDED}. Let $b$ be the amount held by the contract,
$p$ the payer's modeled external holding and $r$ the recipient's holding. All
amounts are nonnegative integer units. Let $d$ be the agreed deposit, $\tau$ the
expiry time and $t$ the current event time. The parties and the meaning of
accepted delivery are agreed inputs to this model.

The state transition table is the executable meaning of the miniature contract.
It says which caller can perform which action, under what guard, and how the
state and resources change. It is deliberately written in a neutral notation;
it is not presented as tested Stipula source code.

\begin{longtable}{@{}P{24mm}P{33mm}P{42mm}P{48mm}@{}}
\caption{Synthetic escrow transitions, inspired by state-contract languages.}\\
\toprule
From & Action and caller & Guard & Effect and next state\\\midrule\endfirsthead
\toprule From & Action and caller & Guard & Effect and next state\\\midrule\endhead
\texttt{UNFUNDED} & Deposit by payer & $p\geq d$, $d>0$, $t<\tau$ & $p'=p-d$, $b'=d$, $r'=r$; enter \texttt{HELD}.\\
\texttt{HELD} & Release by authorized acceptance process & Accepted delivery and $t<\tau$ & $r'=r+b$, $b'=0$, $p'=p$; enter \texttt{PAID}.\\
\texttt{HELD} & Expiry event & $t\geq\tau$ & $p'=p+b$, $b'=0$, $r'=r$; enter \texttt{REFUNDED}.\\
\texttt{PAID} or \texttt{REFUNDED} & Any repeat transfer & No enabled transition & No second transfer is permitted by this model.\\
\bottomrule
\end{longtable}

Start with $p=100$, $b=0$, $r=0$, and an agreed deposit $d=40$. A valid deposit
produces $(p,b,r)=(60,40,0)$ and state \texttt{HELD}. If delivery is accepted before
expiry, release produces $(60,0,40)$ and state \texttt{PAID}. If expiry occurs
first, the refund produces $(100,0,0)$ and state \texttt{REFUNDED}. The same deposit
has two possible future paths, determined by events and their timing.

A reader can now check a conservation property without trusting a compiler. Define
$W=p+b+r$, the total modeled resource. For the deposit transition,
\begin{equation}
 W'=(p-d)+d+r=p+r=W,
 \label{eq:escrow-deposit}
\end{equation}
because the initial contract holding is zero. For release,
$W'=p+0+(r+b)=W$. For refund, $W'=(p+b)+0+r=W$. Each permitted transition therefore
preserves the total. Since the initial holdings are nonnegative and the deposit
guard requires sufficient funds, these transitions also preserve nonnegativity.
This is a local derivation about the stated miniature model.

The terminal states prevent a second transfer. After release, no transition from
\texttt{PAID} can release the same balance again. A Java implementation that
updates balances but forgets the state change could violate that property under
a repeated call. This is why the state transition is part of the contract, not
just a user-interface label. An idempotency mechanism can support repeated
requests, but must be specified consistently with the state model.

% BEGIN REVISION ADDITION teach-M04-06
\begin{figure}[H]
\centering
\TeachingBranch{Deposit held in escrow}{Successful condition: transfer}{Alternative condition: refund}
\caption{The escrow example separates authorized paths through the same resource state.}\label{fig:teach-M04-06}
\end{figure}

For a synthetic deposit of 100 units, a complete refund and a complete transfer cannot both spend the same held deposit. Follow one transition at a time and check the remaining balance. This illustrates the resource argument; the local contract's permitted transitions determine when either path can occur.


% END REVISION ADDITION teach-M04-06
Time is equally substantive. The table uses $t<\tau$ for release and
$t\geq\tau$ for expiry. At exactly $\tau$, only expiry is enabled. If the intended
agreement allowed delivery at the deadline, the guards would need a different
boundary or an explicit ordering rule. A test performed one minute before expiry
cannot establish the exact-deadline behavior. The Stipula literature's event
ordering is therefore a modeling choice to understand, not a universal answer
for every contract \citep{stipula2021}.

The model also exposes a dependency hidden by ordinary prose: what establishes
accepted delivery? It might be an authorized person's recorded acknowledgment,
an external system event or a disputed factual assessment. The transition table
assumes that this predicate is supplied correctly. Conservation and no-double-
transfer proofs do not establish that assumption. If acceptance can be forged,
the formal transfer properties may hold while the wrong person receives the
resource.

The distinction between safety and progress is now visible. Conservation is a
safety property: no permitted step creates or destroys the modeled resource.
The ability to reach \texttt{PAID} or \texttt{REFUNDED} is a progress question.
If the platform never processes the expiry event and delivery is never accepted,
the resource can remain held. Proving eventual refund requires assumptions about
time advancement and event execution, not merely the transition table's arithmetic.
The contract-liquidity literature concerns such availability of actions and
resources; it should not be confused with market liquidity \citep{liquidity2023}.

This miniature also clarifies what a Java/JML translation would need to preserve.
A precondition specifies the caller, state and guard. A postcondition specifies
the new balances and state. An invariant states nonnegative holdings and any
conserved total. The verifier reasons about all allowed calls under those
conditions. It must also account for the event mechanism and any loop abstraction
used by the translation. The inspected Stipula-to-Java work provides a relevant
route, with the restrictions already described \citep{stipulakey2025}.

For LegalMath, the corresponding benefit is a disciplined way to model review,
consent and obligation lifecycles. A report awaiting review is not an approved
report; a withdrawn consent is not an active generation; an opened duty is not
fulfilled merely because a document was eventually delivered late. Explicit
states and transitions make these distinctions executable and testable.

The correspondence is methodological rather than a claim that every circular is
an escrow contract. Regulatory obligations can apply even when a party never
agreed to a private contract, and actual events must be recorded even if a formal
platform would reject them. A bank's monitoring system therefore separates the
normative question of what should occur from the factual record of what did
occur. The state model helps express that separation only when it is designed
for the intended purpose.

A practical specification exercise follows directly. Before implementing a new
history-dependent control, write its states, actors, events, guards and effects;
execute one normal path, one boundary path and one violating observation by hand;
identify the facts that the model assumes; and state which invariants and progress
properties matter. Then choose a language or verifier that supports those
properties. This sequence gives Stipula, Symboleo, eFLINT or a restricted custom
monitor a concrete role rather than selecting a language because its name occurs
frequently in the literature.
% END SOURCE UNIT M04:06

% BEGIN SOURCE UNIT M04:07
\section{What contract verification does and does not buy}
\label{merge:M04:07}

\label{sec:languages-contractverification}
Once a contract has a transition semantics, tools can ask whether certain
states are reachable, whether a resource can become trapped, or whether two
contracts have corresponding observable behavior. Stipula's legal-bisimulation
work gives a formal account of selected observable equivalences
\citep[behavioral equivalence section]{stipula2021}. A bisimulation relates
states so that an observable step on one side can be matched by a corresponding
step on the other. It is a statement about modeled behavior under that
definition of observation.

For the bank, an analogous question might ask whether a revised consent
implementation preserves the reviewed behavior for grants, withdrawals and
transactions. The observation set must include what matters: subject, category,
time, permitted use and recorded violations. If the comparison ignores
withdrawal, equivalence of the remaining observations would not establish the
property the bank needs. Formal observation boundaries deserve the same
scrutiny as ordinary input-domain restrictions.

\citet{stipulakey2025} demonstrate a translation from Stipula contracts into
Java with JML specifications and verification using KeY. JML expresses properties
of Java behavior, such as preconditions, postconditions and invariants; KeY
reasons about whether the Java program satisfies them. This is especially
relevant to our Java delivery target because it illustrates a route from a
specialized contract language into a general-purpose verification environment.

% BEGIN REVISION ADDITION teach-M04-07
\begin{figure}[H]
\centering
\TeachingFlow{Stated contract property}{Verification under its restrictions}{Result for the modeled contract}
\caption{A contract proof inherits the model's expressiveness and assumptions.}\label{fig:teach-M04-07}
\end{figure}

% END REVISION ADDITION teach-M04-07
The inspected work has explicit restrictions. Its event conditions are modeled
symbolically, and the automated encoding keeps certain loop inputs constant
across iterations. Its examples establish feasibility for the examined models,
not complete fidelity to arbitrary calendars or external event streams. The
official translator source was inspected in the earlier project review; no KeY
proof was replayed. Adopting a similar route would require a precise mapping
from our RuleIR and event profile to the Java/JML semantics, including the
unknown and conflict results that ordinary Boolean contracts omit.

General verification cannot be promised for every expressive contract language.
\citet{reachability2026} establish undecidability results for clause reachability
in Stipula fragments and identify a decidable fragment under particular
restrictions. The restrictions concern disjoint function/event source states
and immediate events; they should not be confused with the different
disjoint-cycle condition used in the Java/JML translation work. A shared word
in two papers does not identify a shared theorem.

The practical consequence is to define bounded questions and supported
fragments. A tool can report that no bad state was found up to a specified
depth, or prove a property under a restricted transition model. It cannot
turn a finite search into a complete reachability decision for a language
whose general problem is undecidable. The report must retain the bounds so a
reviewer can understand which behaviors were left unexplored.

Stipula's amendment literature is also relevant. A contract can change its
declarations, state and permissible behavior through a specified amendment
mechanism \citep[secs.~3--5]{amending2023}. For regulatory controls, an amendment
must similarly identify whether existing obligations keep their opening source
version or migrate to a new one. Hot-replacing code without a state-migration
interpretation can change the treatment of historical events. The versioned
control and its transition obligations need separate approval.
% END SOURCE UNIT M04:07

% BEGIN SOURCE UNIT P-01a-language-lesson:04
\subsection{From a consent property to a Java verification model}
\label{merge:P-01a-language-lesson:04}


A property is a precise question about the model. For the consent example:
``After a known withdrawal, and before a new qualifying agreement, this control
never returns that the streamlining conditions are met.'' Testing checks selected
histories. Model checking explores all histories within a specified finite
model. A mathematical proof establishes the property under explicit assumptions.
None establishes that a real client signed a document, or that the selected
provisions are the complete legal regime.

Stipula's behavioral-equivalence work compares the observable interactions of
two formal contracts. Two implementations can use different internal state names
yet be equivalent if the relevant party actions, communications and asset
behavior match. This is useful for comparing formal implementations. It is not
a test that the English circular and the contract mean the same thing
\citep[sec.~5]{stipula2021}.

% BEGIN REVISION ADDITION teach-P-01a-language-lesson-04
\begin{figure}[H]
\centering
\TeachingFlow{Consent-state property}{Java model and assumptions}{Verification obligation}
\caption{A Java verification model must preserve the event meaning it is meant to check.}\label{fig:teach-P-01a-language-lesson-04}
\end{figure}

% END REVISION ADDITION teach-P-01a-language-lesson-04
The Stipula--KeY work supplies a particularly relevant connection to Java.
Its translator builds Java verification models from a supported class of
Stipula contracts and adds specifications in the Java Modeling Language (JML).
KeY reasons about Java programs and these specifications. A JML precondition
states what must hold before a method; a postcondition states what the method
must establish; an invariant states a condition preserved at the designated
program points. For instance, an illustrative postcondition can require a
withdrawal method to leave consent inactive, while an invariant can prevent a
successful streamlining decision in that state. The proof obligation concerns
all executions covered by the model and its assumptions
\citep[secs.~3--5]{stipulakey2025,stipulakeycode}.

That translation is a useful research route, but its Java verification model is
not automatically a deployable bank integration. The paper restricts cycle
structure and abstracts some event conditions; the adapter still has to preserve
timestamps, evidence and the bank's actual interface. The concrete Java
demonstration in this proposal instead compiles the stated decision subset
directly. It compares Java with a separately implemented reference evaluator
and source-derived examples. A later KeY or Lean task can strengthen the
compiler evidence for a chosen property without changing who approves the
legal interpretation.
% END SOURCE UNIT P-01a-language-lesson:04

% BEGIN SOURCE UNIT P-02-literature:05
\subsection{The evidence and restrictions behind Stipula verification}
\label{merge:P-02-literature:05}


Stipula makes parties, agreement, state, assets, permitted functions and scheduled
events explicit. Its semantics describes how a contract evolves when parties act
and time advances. This is well aligned with consent, withdrawals, review events
and requests for evidence. It also provides a notion of legal bisimulation,
relating contracts through their observable interactions
\citep[secs.~3--5]{stipula2021}.

The important borrowing is the state model, not a presumption that bank compliance
requires a blockchain. The original paper discusses centralized as well as
distributed execution. A bank can record a source-linked event history in its
ordinary systems. Moreover, a similarity of formal interaction does not establish
equivalence of the source texts as a matter of law. The accompanying discussion
of legal calculi highlights how executable arrangements interact with external
intervention and legal choices that code can otherwise suppress
\citep{stipula2021,crafa2022}.

Several semantic choices need explicit review. In the studied Stipula semantics,
scheduled events take precedence over function execution at the same time. That
is a clear rule for the language; it is not an automatic answer to the bank's
cut-off or consent-ordering questions. The prototype should construct a trace in
which a withdrawal and a transaction have the same recorded time, and identify
the finer ordering evidence or adopted policy required to decide it.

% BEGIN REVISION ADDITION teach-P-02-literature-05
\begin{figure}[H]
\centering
\TeachingCompare{Published theorem}{Tool's supported fragment}{Bank workflow to be modeled}
\caption{The theorem's restrictions must survive the proposed transfer.}\label{fig:teach-P-02-literature-05}
\end{figure}

% END REVISION ADDITION teach-P-02-literature-05
The Stipula-to-Java/JML work offers a concrete verification route. It translates
contract structure into Java together with specifications and applies KeY to
properties of the result. The authors report automatic verification of four contract examples for the
selected translation and properties. This is reported feasibility evidence
whose exact executable inputs must still be reconciled with the appendix. The inspected implementation includes a
translator and generated examples \citep[secs.~3--4 and apps.~A--B]{stipulakey2025,stipulakeycode}.
The retained extended version also contains a concrete reproducibility
problem. Appendix B.3 declares Boolean parameters in methods that add
those parameters to integer asset balances. Java rejects that arithmetic;
the printed example is not compilable as written. A rendering check
confirmed the declarations, and a minimal compiler diagnostic confirmed
the type error. The pinned official repository also contains the same
Boolean parameters in \texttt{Deposit.java}; compiling that complete
file produces eleven Java type errors \citep{stipulakeycode}.
This does not determine what files were used for the
reported experiment. Adoption requires the exact verified sources,
correction of the discrepancy and a replay of the proof obligations
\citep[app.~B.3, pp.~31--32]{stipulakey2025}.

This is worth a bounded experiment if workflow verification proves important.

The encoding also defines the limit of the evidence. Event conditions are
abstracted through symbolic Boolean values rather than a complete implementation
of banking calendars. Automated handling of the relevant loops uses inputs that
remain constant across iterations. A proof about that model does not establish
correctness for arbitrary changing portfolios, unbounded event queues or all
real-time orderings. LegalMath must state and test any adapter's correspondence
to its own event semantics before using the proof in a compliance package.

The 2026 reachability paper supplies a direct design constraint. Through
counter-machine reductions it establishes undecidability for general and several
restricted microStipula fragments. It also identifies a decidable fragment,
called DI, with both disjoint function/event source states and immediate events,
using a well-structured transition-system argument
\citep[secs.~2--4]{reachability2026}. DI is not the same condition as the
disjoint-cycle construction in the Java/JML paper.

The product implication is to ask bounded, named questions. For example, can an
approved rule set permit streamlined treatment after a recorded withdrawal in
any trace of at most $k$ events under a specified ordering model? A counterexample
is useful immediately. Failure to find one supports only the checked bound unless
an additional invariant proof establishes a broader claim. A timeout or unknown
result must remain unresolved. The research does not imply that a small acyclic
financial decision is intractable; it limits the promise of a universal verifier
for unrestricted event-based contracts.
% END SOURCE UNIT P-02-literature:05

% BEGIN SOURCE UNIT P-02b-contract-evolution:00
\subsection{Amendment and resource analysis}
\label{merge:P-02b-contract-evolution:00}


The amendment literature asks what happens to a running agreement when its code
changes. Higher-order Stipula lets an amendment remove declarations, add
declarations and run a transition, changing both program and current state.
Restrictions can preserve parties, constant fields or asset conditions; an
agreement predicate can require relevant parties to accept the change
\citep[secs.~4--6]{amending2023}. The retained source is the 2023 author manuscript
\emph{Legal Contracts Amending with Stipula}. It is not mislabeled as the different
2024 chapter \emph{Programming Contract Amending}, whose full text was not
successfully retrieved.

% BEGIN REVISION ADDITION teach-P-02b-contract-evolution-00
\begin{figure}[H]
\centering
\TeachingFlow{Original contract state}{Adopted amendment relation}{Resource and duty consequences}
\caption{Amendment changes more than the text of a condition.}\label{fig:teach-P-02b-contract-evolution-00}
\end{figure}

% END REVISION ADDITION teach-P-02b-contract-evolution-00
This poses a precise SFC transfer question. A new circular can change future
decisions while existing consent arrangements or open duties need a separately
justified transition. A replacement marker alone cannot decide that migration.
The MVP creates an immutable new rule version and retains the version that
opened each obligation. Any migration is an explicit reviewed operation.
Runtime code replacement is unnecessary for the initial prototype and would make
historical replay harder to explain.

The platform paper describes a visual editor, interpreter, type inference,
unreachable-clause analysis and liquidity analysis
\citep[secs.~3--7]{stipulaplatform2026}. Its liquidity property concerns resources
that might remain trapped in a program. The underlying type analysis represents
asset effects symbolically; cheaper sufficient tests differ from more precise
bounded exploration \citep[secs.~4--5 and proof appendices]{liquidity2023}.
This is program-resource liquidity, not market or funding liquidity. The platform
account assumes cooperative, strategyless behavior for the discussed property
and states that the liquidity implementation is still being developed.
That is the paper's historical status: the separately retained workbench
repository includes a Python liquidity analyzer and a configurable bound
on function repetitions \citep{stipulacode}. Source inspection confirms
its presence, but neither its correctness nor production suitability has
been established by execution here. Borrowing
editor and explanation ideas is therefore a different commitment from claiming
a production verifier for all bank workflows.
% END SOURCE UNIT P-02b-contract-evolution:00

% BEGIN SOURCE UNIT M04:08
\section{Other normative languages clarify different distinctions}
\label{merge:M04:08}

\label{sec:languages-other}
Symboleo models legal contracts using concepts such as obligations and powers,
with lifecycle distinctions that are easy to lose in a Boolean implementation
\citep[secs.~II--IV]{symboleo2020}. Creating a duty is not always the same as
activating it. A power can change a normative relationship, for example by
suspending or terminating an obligation. These distinctions are useful when a
bank workflow has review, activation, withdrawal and corrective actions with
different legal effects.

The retained Symboleo paper and software inspection do not provide a complete
implemented axiom set for LegalMath. Some supporting references in the early
paper are not fully available in the retained edition. We can use the lifecycle
questions to improve requirements, while treating any direct engine adoption
as a separate interface and semantics audit. A familiar conceptual diagram is
not a substitute for an executable contract.

eFLINT separates facts, actions and duties in a normative specification
\citep[secs.~3--5]{eflint2020}. A particularly useful distinction is that an
impermissible action may still occur, change state and create a violation.
This matches the difference between a bank's pre-transaction advice and its
post-event audit. A system that refuses to represent an unauthorized observed
action cannot explain what happened after a control failure.

% BEGIN REVISION ADDITION teach-M04-08
\begin{figure}[H]
\centering
\TeachingCompare{Permission to act}{Observed action}{Consequences of violation}
\caption{A normative model can record an action that was not permitted.}\label{fig:teach-M04-08}
\end{figure}

% END REVISION ADDITION teach-M04-08
The eFLINT literature also shows why language versions matter. The later
stable-model semantics provides a particular formal grounding and can produce
no model or several models \citep[translation and evaluation sections]{eflintstable2025}.
Those outcomes carry information; selecting one model silently would conceal
ambiguity or inconsistency. The newer language survey discusses implementations
whose semantic guarantees differ \citep{eflint2026}. A proof or test result for
one version should not be attributed to every interpreter bearing the same name.

s(LAW) combines legal formalization with logic programming that can explain
positive and negative outcomes and explore possible models
\citep[secs.~2--5]{slaw2024}. This makes assumptions visible in a way that a
single opaque classification may not. However, a model containing a hypothetical
fact does not establish that the fact is true for the client. Counterfactual
reasoning should propose what would change an answer, while the evidence
service establishes whether the corresponding condition actually holds.

These methods suggest an interface in which normative classification,
state transition and factual evidence are separate inputs to a decision. The
implementation should reject an unsupported duty kind explicitly. A maintenance
obligation, which must hold throughout an interval, cannot be implemented by
checking that it was true once. A non-preemptive achievement obligation cannot
be discharged by a performance before it came into force. The taxonomy in
\citet[sec.~3]{normative2016} helps identify these choices; the precise profile
still needs a source-grounded interpretation.
% END SOURCE UNIT M04:08

% BEGIN SOURCE UNIT P-02-literature:04
\subsection{Authority, default negation and competing normative backends}
\label{merge:P-02-literature:04}


LegalRuleML is relevant because it treats legal meaning as more than a Boolean
expression. The standard addresses modalities, defeasibility, temporal aspects,
sources and contextual information. Earlier work on legal interpretations shows
how alternative readings and their contexts can be represented rather than
collapsed prematurely into one rule
\citep[sec.~4]{legalinterpretations2014}. The standard describes the relevant
requirements and contextual annotations in sections 2.2--2.3
\citep{legalruleml2021}. This is a strong basis
for LegalMath's metadata: every executable condition should identify its source,
scope, adopted interpretation and relationship to other rules.

The standard is an interchange representation, however. XML that names an
obligation does not by itself define the event processor that creates, satisfies
or breaches it. LegalMath should adopt a deliberately small execution profile
with explicit meaning, and use LegalRuleML concepts for interoperability. Requiring
compliance officers to review raw XML would not solve the bank's usability problem.

s(LAW) illustrates a different benefit of formal legal reasoning: explaining
conclusions under alternative factual assumptions. Its use of s(CASP) combines
logic-programming rules, explicit negative information and default negation, and
produces justifications. The school-admission examples demonstrate discretionary
conditions and hypothetical cases, including assumptions that affect which
conclusions can be derived \citep[secs.~4--5]{slaw2024}.

This distinction matters when an SPI record lacks an agreement or a qualification.
A hypothetical answer can say which additional fact would establish a condition.
It cannot treat that fact as already established. LegalMath should display
``would hold if this evidence were confirmed'' separately from an evaluated
client fact. Different treatments of negation must also remain explicit: absence
from a database, a documented negative statement and a defeasible assumption are
three different things. The s(LAW) examples support explanation and exploration;
the linked implementation could not be fully audited in this review.

eFLINT contributes an especially useful operational distinction. A prohibited or
otherwise impermissible action can still occur in the real world and have legal
effects; the system must be able to record the resulting violation. A pure
authorization engine that simply removes such transitions cannot audit what
actually happened. The language's duties, actions and open types are relevant
to missing evidence and ongoing controls \citep[secs.~3--5]{eflint2026}.

The 2026 eFLINT account also describes differences between implementations and
versions. It notes the absence of a formal operational semantics and type system
for version 4, and limitations concerning fine-grained explanations and normative
priority; another implementation has a stable-model foundation. Guarantees must
therefore be attributed to the actual version and engine used
\citep[secs.~6--7]{eflint2026}. The prototype should borrow the distinction between
permission, occurrence and violation without assuming that installing an eFLINT
interpreter supplies all the required proof evidence.

% BEGIN REVISION ADDITION teach-P-02-literature-04
\begin{figure}[H]
\centering
\TeachingCompare{No proof of a fact}{Explicit negative fact}{Conflict between authorities}
\caption{Different normative backends distinguish absence and negation differently.}\label{fig:teach-P-02-literature-04}
\end{figure}

% END REVISION ADDITION teach-P-02-literature-04
The original eFLINT paper makes the operational idea concrete through a
transition-system interface: declarations establish institutional concepts,
events change a configuration, and queries inspect permissions and duties.
Action-compliance and duty-compliance are independent. Its GDPR example forms
part of a small data-sharing/KYC exploration with ABN AMRO and ING, using separate
policy and agreement specifications. This is relevant domain evidence, although
it is a selective experimental case rather than a deployed private-bank control
system \citep[secs.~3--5]{eflint2020}. Operational connectors qualify observations
into facts; the normative engine evaluates those facts; a reviewer owns the
interpretation connecting them.

The stable-model implementation supplies a sharper warning about backend
substitution. \citet[secs.~4--7]{eflintstable2025} encode states and transitions
in Clingo and translate a restricted normal form of eFLINT clauses. The
translation makes state explicit, normalizes aggregation and permits scenario
search as well as checking a supplied scenario. Section 8 compares 108
specification--scenario pairs over 43 specifications. Agreement with the earlier
interpreter is expected only where their semantics coincide; default reasoning
is deliberately changed.

Their Example 9 shows how a prior implementation can retain a conclusion based
on absence even after further facts invalidate that absence. Examples 10 and 11
distinguish no stable model from multiple stable models, while the older engine
produces a particular result. A future ASP adapter that selects the first model
silently would therefore choose an interpretation policy. LegalMath 0.1 avoids
that problem by rejecting recursive definitions and defining unknowns directly.
Richer semantics remain an option when the selected regulatory scope requires
them. Pinning a package name alone does not pin the meaning of a rule.

To understand a stable model, consider the two rules ``accept if rejection is
not established'' and ``reject if acceptance is not established''. One
self-consistent interpretation contains acceptance; another contains rejection.
Selecting the first answer chooses between them without additional legal
authority. In answer-set programming (ASP), a solver enumerates such stable
interpretations after accounting for the program's default-negation rules.
Absence of a stable model can indicate inconsistent constraints, while several models can leave
the answer dependent on the chosen interpretation. This miniature example
explains the backend issue without equating stable-model semantics with the
three-valued fact evaluation used by LegalMath.

The accompanying eFLINT release makes these limits inspectable: it separates
tests expected to agree from tests expected to differ, and documents excluded
cases involving cyclic aggregation and unsatisfied models. The compiler's
\texttt{Clingo.hs} also keeps action violations separate from duty violations.
The source was inspected, not compiled or benchmarked. This supports a small
differential test before adopting an engine, rather than extrapolation from the
paper's overall claims \citep{eflintartifact}.

Symboleo contributes an explicit lifecycle for obligations and legal powers.
A trigger creates a conditional obligation; its antecedent activates it; later
events can fulfill, violate, suspend or terminate it. A power concerns the ability
to change these legal positions, which differs from a duty to act. The initial
paper gives statecharts, temporal predicates and selected axioms, with a
sale-of-goods example checked using a Prolog prototype
\citep[secs.~III--IV]{symboleo2020}. It is preliminary feasibility evidence;
the retrieved preprint does not contain all 27 referenced axioms. The official
editor and model-checker repositories provide concrete examples and a nuXmv
route, but were not executed \citep{symboleocode}. The borrowing for LegalMath
is to distinguish obligation creation, activation and performance. Suspension
and legal powers require an extension to the MVP's simpler profile.
% END SOURCE UNIT P-02-literature:04

% BEGIN SOURCE UNIT M04:09
\section{Decision tables and business-rule engines}
\label{merge:M04:09}

\label{sec:languages-tables}
Decision tables remain useful when their rows and hit policy have a precise
meaning. A table can show, for example, that a known discount exception changes
the selected trigger while an unknown exception leaves a material question.
The reader must know what happens when several rows match: choose the first,
require a unique row, collect results or apply a priority. DMN makes such hit
policies explicit \citep[sec.~8.2.11]{dmn2024}.

That explicitness is the valuable part. A spreadsheet whose row order silently
resolves competing legal exceptions is harder to review because its legal
priority is hidden in presentation. Changing the sort order can change the
answer. A unique-hit table can reject overlap, but it cannot determine whether
the missing row represents an omitted legal situation. Independent source
coverage and boundary cases remain necessary.

% BEGIN REVISION ADDITION teach-M04-09
\begin{figure}[H]
\centering
\TeachingFlow{Decision-table conditions}{Declared hit policy}{One specified result}
\caption{A table needs an explicit rule for multiple matching rows.}\label{fig:teach-M04-09}
\end{figure}

If two rows match the same client, a first-match policy and a collect policy compute different outputs. Visual neatness does not choose between them. The bank must adopt the intended meaning before the table becomes an interchangeable backend.


% END REVISION ADDITION teach-M04-09
Drools, OPA and other rule engines can execute policy logic and integrate with
enterprise applications. The question for this product is not whether they
are capable software. It is whether the selected engine's semantics, evidence
model, explanation behavior and deployment interface match the adopted
specification. Unknowns, conflicts, exception priorities and event ordering
need explicit mappings. A backend name is not an assurance claim.

For the prototype, a small RuleIR is easier to specify and compare than a
large engine with many unsupported features. Other engines can serve as
independent reference encodings for a supported fragment after that fragment
is defined. The comparison must include the behavior that ordinary Boolean
examples leave out. Engine diversity without semantic alignment can produce
many discrepancies that reflect adapter choices rather than flaws in the
legal interpretation.
% END SOURCE UNIT M04:09

% BEGIN SOURCE UNIT P-02-literature:08
\subsection{Software adoption and interchange standards}
\label{merge:P-02-literature:08}


The open-source landscape is sufficiently rich that a prototype need not build
every component. It is not sufficiently uniform that selecting one package
settles semantics, review, execution and proof at once. Table~\ref{tab:software}
identifies a practical role for each candidate. These are design recommendations
based on the inspected papers and official documentation; installation,
performance, security and integration have not been tested here.

The candidates fall into four useful families. The point of the comparison is
not to choose the most famous package. It is to ask which distinction the package
can preserve, and which part of the bank's problem would still require a separate
record or review. The compact table gives that map; Appendix~\ref{app:detailed-tables}
retains the candidate-by-candidate dispositions and citations.

\begin{table}[htbp]
\centering
\caption{What the main language families can contribute.}
\label{tab:software}
\begin{tabularx}{\textwidth}{P{35mm}P{52mm}Y}
\toprule
Family & Useful contribution & Boundary to keep visible\\
\midrule
Source-linked rule languages & Definitions, exceptions and source-linked authoring, as in Catala. & A readable rule still needs an event history, evidence policy and independent review.\\
Normative and contract languages & Actors, duties, permissions, states and transitions, as in LegalRuleML, Stipula, eFLINT and Symboleo. & Their models do not by themselves settle the bank's source meaning or data mappings.\\
Decision and policy engines & Executable tables and services, as in DMN, Drools, L4 and Open Policy Agent. & Hit policies, missing values, explanations and legal provenance need explicit alignment.\\
Solvers and research tools & Counterexamples, bounded search and selected proof obligations, including Z3 and Lean. & A solver result is evidence about a specified rule; it is not evidence that the rule was legally chosen.\\
\bottomrule
\end{tabularx}
\end{table}

% BEGIN REVISION ADDITION teach-P-02-literature-08
\begin{figure}[H]
\centering
\TeachingCompare{Open-source implementation}{Interchange specification}{Approved bank dependency}
\caption{Availability, interoperability and adoption are separate questions.}\label{fig:teach-P-02-literature-08}
\end{figure}

% END REVISION ADDITION teach-P-02-literature-08
L4 is a serious alternative to building the authoring environment from scratch.
Its inspected repository exposes a parser, type checker, evaluator, IDE,
decision service, generated schemas and graphical traces. Its regulative
\texttt{DEONTIC} documentation names actor, action, deadline, successful
continuation and breach continuation, and distinguishes \texttt{MUST},
\texttt{MAY} and \texttt{SHANT} \citep{l4code}. The 2023 L4 paper on deontics
and time was identified but its full text was not retrieved, so no theorem from
it is used here. Before adopting L4 as a core dependency, test exact money,
missing evidence, simultaneous events and historical replay against the RuleIR
cases. The proposal fixes that comparison target rather than making development
wait for another language-selection exercise.

DMN deserves particular attention because its accessible tables can conceal a
priority choice. Under a Unique hit policy, matching rows must be disjoint; Any
permits overlaps with the same result; First selects by row order. Those policies
are not interchangeable. A default exception tree imported from Catala cannot be
exported as an arbitrary First table without checking that the meanings agree
\citep[sec.~8.2.11]{dmn2024}. Review annotations also need their own source store:
a visual note on a table is not automatically available as runtime provenance.

Licensing and reproducibility are separate from public availability. Catala,
Apache KIE, OpenFisca, Z3 and the other public projects are candidates for reuse,
but the precise pinned version, dependencies and intended distribution must be
reviewed. A paper PDF, a hosted service and a public research repository carry
different rights and maintenance expectations. The Stipula and tax prototypes
should not be described as approved bank components merely because their source
can be inspected. No third-party implementation has been imported into LegalMath
as part of this document work.
% END SOURCE UNIT P-02-literature:08

% BEGIN SOURCE UNIT M04:10
\section{A solver can improve reasoning after a translation}
\label{merge:M04:10}

\label{sec:languages-neurosymbolic}
A neural-symbolic system separates tasks that language models and formal
engines perform differently. A language model proposes a formal representation
of prose. A symbolic engine evaluates or reasons about that representation.
The engine can give a reproducible answer to a precise logical question and
can expose inconsistency, unsupported syntax or a counterexample. The
translation that supplied its premises remains a separate source of error.

Logic-LM uses language-model formulation followed by an appropriate solver and
feedback-driven refinement \citep[sec.~3 and prompt appendices]{logiclm2023}.
LINC combines language-model formalization with first-order theorem proving
and analyzes failures at the translation and reasoning boundary
\citep[method and error analysis]{linc2023}. These methods support the design
choice that a model should not be asked to simulate every logical calculation
in prose. They do not support treating successful solver execution as evidence
that a legal requirement was fully captured.

A syntax repair illustrates the benefit and the limit. Suppose a candidate
uses an undeclared symbol in the gift rule. The parser identifies the symbol
and rejects the encoding. A repair prompt can ask for the missing declaration
or corrected reference while preserving the intended interpretation. The
repaired candidate is then validated again. This process may solve a concrete
representation error. It does not justify adding a new legal assumption
merely because doing so makes the formula easier to satisfy.

ARc develops policy-model creation, human vetting, structured inspection and
symbolic checking in a combined system \citep[sec.~3 and appendix A]{arc2025}.
Its policy creation procedure divides a document into units, proposes types,
variables and constraints, and composes them into a model. Its vetting
mechanisms include linting, template-based English inspection and testing.
These are relevant to our product because they make the proposed model
available for correction before individual answers rely on it.

% BEGIN REVISION ADDITION teach-M04-10
\begin{figure}[H]
\centering
\TeachingFlow{Language-model formalization}{Solver evaluates that formula}{Answer conditional on translation}
\caption{A solver can reason correctly about an incorrectly translated question.}\label{fig:teach-M04-10}
\end{figure}

% END REVISION ADDITION teach-M04-10
Composition creates its own interpretation problem. Two units may use the same
word for different legal concepts, or different words for the same concept.
ARc uses embedding-based clustering in its variable-unification procedure.
For our high-stakes specification, such a merge should be a proposal with
provenance and validation. Merging ``offer discount'' with ``component
discount'' because their descriptions are similar could remove the distinction
that g04 was designed to test. A similarity score cannot establish identity
of legal concepts.

ARc's answer-verification procedure also uses redundant translations and
symbolic agreement checks. Its reported finding categories distinguish
ambiguous translation, impossible premises, invalid conclusions, entailed
conclusions and satisfiable but unentailed claims. The distinction is useful:
if a model permits both an outcome and its opposite under incomplete premises,
the system can show scenarios for each rather than invent a definitive answer.
For LegalMath, these scenarios become review questions tied to the candidate
set and the source context.

The word ``soundness'' in ARc's empirical evaluation needs careful translation
into a product claim. The paper defines an empirical quantity based on false
positives across requests; it is not a theorem that the English policy was
formalized correctly. The denominator differs from precision among approvals.
The manuscript therefore borrows the inspection and disagreement mechanisms
without importing a claimed legal-error probability. Chapter~\ref{ch:evaluation}
defines the denominators our own evaluation would need.

The tax-reasoning work of \citet{jurayj2026} further illustrates the importance
of task boundaries. Supplying gold rules, retrieving annotated cases and
extracting facts are different assistance conditions from generating rules
from a new circular. A comparison must account for those inputs and the human
effort that produced them. A system given a reviewed rule library should not
be presented as evidence that the same system can independently interpret
unreviewed prose.
% END SOURCE UNIT M04:10

% BEGIN SOURCE UNIT P-02-literature:07
\subsection{Language models can propose formalizations; their inputs remain the weak point}
\label{merge:P-02-literature:07}


The SARA tax-reasoning dataset makes explicit the ingredients an evaluation needs:
statutory text, cases, formal rules and expected conclusions. Its corpus uses
selected and simplified US tax provisions and constructed examples. It is a
valuable experimental design, but neither a current tax authority nor a test set
representative of SFC circulars \citep[secs.~2--3]{sara2020}.

The later work connecting statutory reasoning to information extraction associates
facts with spans in case descriptions and evaluates their effect on reasoning.
This is a useful model for recording where a client attribute came from, which
person it describes and how it was normalized
\citep[secs.~3--5]{connecting2023}. Its task-specific date imputations and
semi-synthetic construction require caution: an extraction convention suitable
for that dataset is not an evidential rule for missing bank records
\citep[sec.~7 and apps.~B, E]{connecting2023}.

% BEGIN REVISION ADDITION gap-nli-evidence
\begin{figure}[H]
\centering
\TeachingCompare{Entailed by the document}{Contradicted by the document}{Not mentioned}
\caption{Evidence-based classification preserves the difference between contradiction and silence.}\label{fig:gap-nli-evidence}
\end{figure}


% END REVISION ADDITION gap-nli-evidence
ContractNLI adds a complementary annotation pattern. It separates supported,
contradicted and not-mentioned hypotheses, with evidence spans in contracts.
Its 607 non-disclosure agreements and fixed hypotheses expose the importance of
exceptions and evidence spread across a document. It supplies useful ideas for
an evidence-review interface, while evaluating document inference rather than
compilation of regulatory rules \citep[secs.~2--5]{contractnli2021}.

Logic-LM and LINC both divide language interpretation from symbolic inference.
Logic-LM routes formalized problems to solvers and uses error feedback to refine
malformed programs. LINC translates into first-order logic and uses a theorem
prover, with multiple samples and voting in its experiments
\citep[sec.~3]{logiclm2023} \citep[secs.~2--5]{linc2023}. The transferable idea is a
clear boundary: the model proposes the statement, and the solver reasons over
that statement. Their benchmarks do not establish legal translation fidelity.
Successful parsing can leave a wrong quantifier untouched, and repeated models
can agree on the same missing exception.

The inspected Logic-LM repository also offers fallback strategies when formal
execution fails, including a language-model answer or random selection for its
benchmark setting. Those are not acceptable defaults for this prototype's
regulatory decisions. A failed formalization should produce an unresolved issue
with a reason. Research code can be valuable while requiring a different
operational policy \citep{libraryreview}.

The ARc paper is especially close to the proposed product. It separates creation
and human vetting of reusable policy models from subsequent verification of
natural-language material. Its workflow uses source passages, variable
identification, repeated translations and deterministic explanations. The
verification algorithm checks consistency and relationships between the policy,
premises and claim, and distinguishes invalidity, validity, satisfiability and
several forms of inability to decide
\citep[sec.~3 and Algorithm~1]{arc2025}.

The method can be understood with three formulas. Let $M$ be the reviewed policy
model, $P$ the premises extracted from a message and $C$ its proposed conclusion.
For example, the message may say that a client has a HK\$40 million qualifying
portfolio and therefore satisfies the financial condition. The verifier first
checks whether $M$ allows the premises at all. If $M\land P$ has no possible
assignment, the premises are impossible under that model; declaring the
conclusion true by vacuous implication would be misleading. With consistent
premises, if $M\land P\land C$ has no assignment, the conclusion is refuted.
If $M\land P\land\neg C$ has no assignment, it is entailed. If both have
assignments, the premises leave the conclusion undecided; the assignments show
which extra facts distinguish the outcomes. For a purported complete SPI
assessment based only on wealth, those missing facts could include objectives
and knowledge. These checks reason within $M$; they do not reconstruct an
omitted circular paragraph.

ARc's redundant-translation step addresses an earlier source of uncertainty.
Several translations of the same message produce candidate premise/conclusion
pairs. Algorithm 1 counts the fraction of translations that support a pair's
implication while not contradicting its premises. It does not merely compare
the strings or require formula identity. A confidence threshold can make the
system abstain and display competing readings. LegalMath can borrow that
review interaction, but should preserve the candidate formulas and a
distinguishing client example. Agreement among generated candidates is evidence
of agreement, not independent proof that the interpretation is correct.

% BEGIN REVISION ADDITION teach-P-02-literature-07
\begin{figure}[H]
\centering
\TeachingBranch{Natural-language task}{Supplied facts and formal resources}{Automatically produced translation}
\caption{The information supplied to a benchmark limits what its success establishes.}\label{fig:teach-P-02-literature-07}
\end{figure}

% END REVISION ADDITION teach-P-02-literature-07
This is a strong precedent for an interpretation workspace with explicit
abstention. It is also a warning about terminology. The paper's reported
``soundness'' metric is $1-\mathrm{FP}/N$, an empirical rate over its decisions;
it is not a theorem establishing the translator's semantic soundness. At the
3/3 threshold in Table 1, 10 false positives among 1,689 decisions yield about
99.4 percent on that metric, while reported precision is 94.4 percent and valid
recall is 14.9 percent. The 563 items were each run three times, so the total is
not 1,689 independent policies \citep[sec.~4, Table~1]{arc2025}.

Those results motivate measuring both errors and abstention. A cautious system
may produce few confident errors by deciding very little. Its business value
then depends on whether the unresolved cases become easier to review. The ARc
paper also discusses difficulties with cross-references, tables and implicit
assumptions, all present in the bank's problem
\citep[sec.~5 and app.~A.1]{arc2025}. Its semantic comparison of candidate
translations should not be simplified into a claim that independent generation
proves equivalence. ARc is a published commercial-system comparator, not an
open-source implementation established by this review.

The requested AAAI paper on trustworthy tax reasoning evaluates combinations of
language models, facts and logic programs, including reviewed statutory rules and
examples. It provides a particularly useful distinction for prototype design:
reasoning with gold rules is a different task from generating the rules. The
human work needed to prepare those rules and examples belongs in the comparison
\citep[Tables~2--3]{jurayj2026,jurayj2025extended}. The reported setting uses
100 numerical tax cases; its application-specific penalty for errors and abstention
cannot be imported as a Hong Kong bank's cost model.

In implementation terms, the tax method separates at least three jobs: encode
the statute as a reusable logic program, extract the current case's facts, and
execute the program to obtain the requested amount. Supplying human-reviewed
rules removes the first job from the model's evaluation. Supplying formalized
facts also removes much of the second. The corresponding LegalMath experiment
must therefore hold the source packet and human preparation budget fixed, and
report separately rule errors, fact-mapping errors and execution errors. A Java
calculation can be exact while still using a net-asset input that includes the
home; scoring only the final answer hides which stage needs repair.

The inspected companion code confirms that extraction and Prolog processing are
separate stages and that execution failure can lead to abstention. LegalMath
should adopt that separation and evaluate facts and rules independently. It
should not quote a tax benchmark's best configuration as evidence that the same
method will reduce private-bank compliance errors. The reading notes also retain
a discrepancy between the paper's narrative attribution of a cost result and
the configuration shown in its table; the proposal relies on the methodological
distinction rather than that ranking \citep{libraryreview}.

The recent \emph{Know Your Limits} preprint studies failures in legal reasoning
and autoformalization, including assumptions that change the intended entailment
task. Its examples strengthen the case for testing entity binding, implicit
assumptions and distinctions between legal interpretation and strict entailment
\citep[secs.~3--5]{limits2026}. It remains under-review evidence. Inconsistencies
in the reported dataset denominators require clarification before quantitative
reuse, and the choice of which ordinary or legal assumptions are admissible is
itself a modeling decision. This proposal borrows its failure taxonomy, not a
general error-rate estimate.

LegalBench broadens evaluation beyond a single reasoning dataset. Its tasks
separate several kinds of legal reasoning; the evaluation appendix gives
task-specific procedures. For rule application, legal reviewers check both
correctness and whether an explanation contains the inferences connecting the
facts to the result \citep[sec.~3 and evaluation appendix]{legalbench2023}.
Repeating a circular accurately is not an explanation of its application to a
client. Nor can a benchmark score become a banking acceptance criterion: the
numerical SARA task permits error within ten percent, whereas one cent can change
the intended threshold predicate. The limitations appendix identifies US-law,
language, document-length and ambiguity gaps. LegalMath borrows independent
answer guides and error analysis while constructing exact boundaries,
source-dependency tests and amendment histories for its own task.
% END SOURCE UNIT P-02-literature:07

% BEGIN SOURCE UNIT M04:11
\section{Roundtrip verification and targeted repair}
\label{merge:M04:11}

\label{sec:languages-roundtrip}
Roundtrip verification asks whether a formal representation survives a journey
through another representation. In the method of \citet{roundtrip2026}, source
English is formalized, the formula is translated back into English, and the
reconstructed English is formalized again. A solver compares the first and
last formulas. If they differ, diagnosis identifies a likely translation
stage to repair, and dependent stages are regenerated.

This is attractive for our user interface because the intermediate English
can make a formal candidate easier to inspect. Suppose the original formula
contains both product-link routes but the back-translation mentions only a
specific product. Re-formalizing that sentence may lose the product-type
route. The resulting counterexample exposes a defect in the explanation or
translation chain. A reviewer should not be shown that explanation as if it
faithfully described the formal rule.

The same construction can miss a shared omission. If the first formalization
already omits the product-type route, a faithful back-translation of that
incorrect formula and a faithful re-formalization will agree. The roundtrip
then passes. What survived is the first formula, including its error relative
to the source. The paper explicitly recognizes this limit; our ensemble must
retain a source-grounding path that does not simply inspect the same
back-translation.

% BEGIN REVISION ADDITION teach-M04-11
\begin{figure}[H]
\centering
\TeachingCycle{Source meaning proposed}{Translate back and compare}{Repair the identified discrepancy}
\caption{Roundtrip checks can expose loss, but can also repeat a shared omission.}\label{fig:teach-M04-11}
\end{figure}

% END REVISION ADDITION teach-M04-11
Repair therefore needs a named target. If the formal rule is justified but
the explanation drops a condition, repair the explanation and its dependent
roundtrip outputs. If the original rule omits a source condition, create a
new interpretation candidate and revalidate its downstream artifacts. If the
two formulas use incompatible vocabularies, repair the mapping before claiming
the formulas disagree about one question. A generic ``make these consistent''
instruction does not distinguish those cases.

The stopping rule is equally important. A roundtrip procedure can run for a
bounded number of attempts and still end in disagreement. The product should
report the original and revised formulas, the counterexample, diagnosis,
actions attempted and the remaining issue. It should not keep repairing
until some pair agrees and then discard the path that made the original
source objection visible. The ensemble controller extends this principle
from one translation chain to several methods and issue types.

An explanation generated by deterministic templates provides a useful
additional path. It is constrained by the formal syntax and can be checked
for completeness of its displayed operands. It may be less elegant than
free prose, but it offers a stable view of what the program actually says.
A freely generated explanation can still help readers if it is compared
against that view and does not replace it as the sole account of the rule.
% END SOURCE UNIT M04:11

% BEGIN SOURCE UNIT M04:12
\section{Evidence spans and distant qualifications}
\label{merge:M04:12}

\label{sec:languages-evidence}
An answer supported by a source span is easier to review than an unsupported
assertion, but the span itself may be insufficient. A qualification can appear
in a footnote, definition or distant section. ContractNLI studies document
inference with evidence annotation and highlights challenges involving
references and exceptions \citep[task, annotation and error-analysis sections]{contractnli2021}.
For our product, the relevant lesson is to retain the supporting and limiting
passages needed for the inference, rather than a single convenient quotation.

The distinction between ``not mentioned'' and ``false'' is particularly useful.
If a retained source does not settle the definition of a fee refund, the
absence should remain visible. A retrieval engine returning no matching
paragraph does not establish that the exception excludes refunds. It may have
searched the wrong document, missed a synonym or lacked the referenced source.
The evidence report should distinguish unsuccessful retrieval from reviewed
absence within a specified source set.

% BEGIN REVISION ADDITION teach-M04-12
\begin{figure}[H]
\centering
\TeachingFlow{Operative sentence}{Distant definition or footnote}{Qualified meaning}
\caption{The supporting span can lie beyond the nearest sentence.}\label{fig:teach-M04-12}
\end{figure}

% END REVISION ADDITION teach-M04-12
Source spans also require stable identity. Character offsets depend on the
extraction version, while page numbers may not survive a changed layout.
Retain raw-byte identity, derivative identity, exact span text or hash, and
human-readable location. A reference to paragraph 10 without a source version
can become misleading after an amendment. A matching text hash without its
source context may identify a sentence copied into a secondary summary.

The formalization boundary should distinguish source claims from supplied
facts. A sentence defining a gift exception is part of the normative context.
A bank record establishing that a fee was paid is factual evidence. A
reviewer's classification relating the record to the exception is an
interpretive application. Combining all three into an undifferentiated
``evidence'' list makes it harder to see where disagreement actually arises.

Recent faithfulness work analyzes how language models and formal solvers can
diverge when implicit assumptions enter formalization \citep{limits2026}.
The retained paper also has annotation and denominator questions that prevent
us from transferring a general accuracy figure. Its examples nevertheless
reinforce a practical requirement: an implicit assumption that changes the
answer should become an explicit, challengeable premise. The system must
show where a model supplied information that the source packet did not.
% END SOURCE UNIT M04:12

% BEGIN SOURCE UNIT M04:13
\section{Testing paths, dates and inconsistent constraints}
\label{merge:M04:13}

\label{sec:languages-analysis}
Other methods contribute different checks once a formal model exists.
CUTECat uses concolic execution to explore paths in Catala programs
\citep[secs.~2--5]{cutecat2025}. Concolic execution combines a concrete run
with symbolic constraints describing the path taken. By changing a path
condition, the tool seeks an input that exercises another branch. For our
gift example, this is analogous to moving from a named-product case to the
case where only the product-type route can determine the answer.

Path-directed generation can expose an exception branch that ordinary
examples never visit. It cannot visit a branch that the formal program
omitted entirely. Independent source-derived cases remain necessary. The
reported costs on larger legal examples also show why a prototype needs
explicit resource limits; small-program speed should not be treated as a
promise that every circular can be exhaustively explored.

Date arithmetic is another place where a precise program can embody a wrong
choice. Adding a month to a date near the end of a month may require a
rounding convention. The choice can affect deadlines. \citet{dates2024}
formalize date arithmetic and analyze ambiguity in legal applications. The
transfer to this project is to make calendar rules, boundary conventions and
their source basis explicit. A verified date library does not decide which
rounding convention a particular provision requires.

% BEGIN REVISION ADDITION teach-M04-13
\begin{figure}[H]
\centering
\TeachingCompare{Exact threshold equality}{Leap-day boundary}{Contradictory requirements}
\caption{Boundary cases test different sources of failure.}\label{fig:teach-M04-13}
\end{figure}

% END REVISION ADDITION teach-M04-13
ContractCheck analyzes consistency and executability of structured contract
models \citep[secs.~3--8]{contractcheck2025}. A satisfying assignment shows
that a modeled set of conditions can jointly hold. An unsatisfiable set can
help locate conflicting clauses or assumptions. Neither result proves that
the original legal document has been translated faithfully. The study's
manual translation is an important part of its task boundary.

An unsatisfiable core is a useful diagnostic: a subset of constraints that
already suffices for inconsistency. It is not necessarily the smallest such
subset unless a minimization procedure establishes that property. For a
reviewer, a small and source-linked explanation is valuable, but the product
should label a returned core accurately. It must also distinguish an
encoding inconsistency from legitimate defeasible conflict that the chosen
logic was intended to resolve.

These analyses fit into the ensemble as specialized questions. Does this
representation have a hidden path? Does a date convention alter a boundary
case? Are the supplied assumptions jointly possible? Does a revised rule
preserve the reviewed behavior? Their evidence can challenge a candidate
without claiming to replace the source interpretation process.
% END SOURCE UNIT M04:13

% BEGIN SOURCE UNIT P-02-literature:03
\subsection{Dates and test generation expose defects that ordinary examples miss}
\label{merge:P-02-literature:03}


The date-arithmetic work of \citet{dates2024} demonstrates why even a simple phrase
such as adding a month needs a convention. Calendar months have different
lengths; leap days create boundary cases; sequences of date operations need not
behave like addition of integers. The paper develops explicit rounding choices,
formalizes date semantics and selected properties, and uses static analysis to
detect ambiguities in legal computations. Its application to benefit rules
connects the formalism to realistic source language rather than treating dates
as a generic programming-library problem.

The immediate lesson for annual reviews and three-year transaction windows is
to preserve the chosen calendar operation and its source or interpretation.
``Three years'' must not silently become a fixed number of days. A leap-day case
can expose disagreement between two plausible implementations. The paper's
day-level formalization does not settle intraday questions such as withdrawal
instructions arriving at a trading cut-off, nor does the mechanization make the
whole static-analysis implementation verified
\citep[secs.~2--6]{dates2024}. Those are distinct prototype obligations.

% BEGIN REVISION ADDITION teach-P-02-literature-03
\begin{figure}[H]
\centering
\TeachingFlow{Declared date convention}{Generated boundary case}{Observed discrepancy}
\caption{Date semantics needs explicit calendar assumptions.}\label{fig:teach-P-02-literature-03}
\end{figure}

For a synthetic anniversary beginning on 29 February, a non-leap year forces a policy choice. February 28, March 1 and rejection are distinct outcomes. A library default cannot establish which convention the adopted requirement uses.


% END REVISION ADDITION teach-P-02-literature-03
CUTECat complements this work by generating test cases from Catala programs using
concolic execution: an execution carries concrete values and symbolic conditions,
then a solver seeks inputs that take a different branch. The method specializes
its exploration to default expressions and uses preferences to obtain inputs
that are easier for humans to understand. In the reported evaluation, generated
cases exposed 20 manually validated seeded mutations. A much larger example also
shows that exhaustive-looking exploration can produce a substantial volume of
cases and runtime \citep[secs.~3--5]{cutecat2025}.

In the running example, a concrete test with portfolio HK\$45 million takes
the successful portfolio branch. The symbolic execution retains the condition
$P\geq 40{,}000{,}000$ instead of forgetting how that result was reached.
Negating it asks the solver for $P<40{,}000{,}000$; holding the other financial
route false yields a case that changes the overall answer. To distinguish
inclusive from strict comparison specifically, request $P=40{,}000{,}000$.
Concolic exploration automates this alternation between a real execution and
constraints describing another path. Mutation testing asks a different but
complementary question: whether the examples notice an intentionally wrong
operator. Both mechanisms are understandable and useful without importing
the entire CUTECat execution environment.

LegalMath should borrow branch-directed and mutation-based testing, especially
for inclusive boundaries, alternatives, exceptions and rounding. The result is
evidence about the encoded program. A branch corresponding to an omitted
paragraph does not exist for the test generator to discover. Independent
source-derived scenarios are therefore essential. The official CUTECat artifact
documentation was inspected, but its large execution image was not run; the
prototype should first implement a small set of useful mutations before adopting
that full toolchain \citep{libraryreview}.
% END SOURCE UNIT P-02-literature:03

% BEGIN SOURCE UNIT P-02-literature:06
\subsection{Consistency of clauses and compliance of processes}
\label{merge:P-02-literature:06}


ContractCheck models legal contracts through structured clauses, an ontology and
constraints, and uses automated reasoning to identify consistency and execution
problems. Its share-purchase-agreement case study is particularly relevant to
the presentation of contradictions: reviewers need to know which claims cannot
coexist and which circumstances make an obligation executable
\citep[secs.~3--7]{contractcheck2025}.

Two distinctions are essential when borrowing this approach. First, the study's
input formalization is manual; its reasoning results do not measure automatic
translation fidelity. Second, showing that each claim can be executed under
some circumstances differs from showing that all duties can be satisfied in one
execution, and both differ from showing that every possible execution complies.
LegalMath should label each of those questions separately. If a solver returns
an unsatisfiable core, that identifies conflicting constraints; a smallest or
minimal explanation requires a further minimization procedure rather than an
assumption about the core \citep[secs.~4--5, 8]{contractcheck2025}.

Process-compliance research reinforces the need to connect rules with activities,
actors and the effects of execution. \citet{ghose2007} describe effect accumulation,
constraint checking and repair in business processes. For the bank, the analogue
is the difference between determining that a warning is required and showing
that the relevant warning was delivered before the required event.

\citet[sec.~2]{compliancehard2015} distinguish achievement obligations, which must
be accomplished, maintenance obligations, which must hold across an interval,
and punctual obligations. They also distinguish partial compliance, where some
execution complies, from full compliance, where all executions do. Their
complexity results concern specified process and obligation languages
\citep[secs.~3--4]{compliancehard2015}. The prototype should borrow those questions
and delimit its model rather than import a complexity label for all regulation.
The practical consequence is that a flowchart showing one successful path is
insufficient evidence about the paths that omit a review or mishandle an exception.

% BEGIN REVISION ADDITION teach-P-02-literature-06
\begin{figure}[H]
\centering
\TeachingFlow{Norms and process model}{Specified compliance question}{Result under expressiveness limits}
\caption{Clause consistency and compliance of a process are different analyses.}\label{fig:teach-P-02-literature-06}
\end{figure}

% END REVISION ADDITION teach-P-02-literature-06
\citet[sec.~3, Definitions 10--17]{normative2016} give a more detailed taxonomy.
An obligation may apply at a point or persist across an interval; a persistent
obligation may require an achievement once or maintenance throughout; achievement
may or may not count before activation; and a violated obligation may persist or
have specified compensation. A Boolean \texttt{completed} field cannot distinguish
early action, timely performance, late performance, or a continuing control that
failed briefly and recovered. The paper takes as input a function specifying
which obligations are in force at each trace position. Supplying that function
is a substantive modeling step, not something the taxonomy infers from prose.
The printed definition of maintenance violation also needs correction:
Definition 14 on page 15 uses membership in the set of fulfilled conditions
where the adjacent prose and Figure 3.5 require non-membership. As printed,
a trace that satisfies the duty throughout could be called a violation.
An implementation must check the intended condition explicitly instead of
copying that formula.
The MVP consequently supports one named profile, non-preemptive achievement,
and explicitly rejects unsupported maintenance and compensation templates.

\citet[secs.~3--4]{semanticcompliance2016} connect LegalRuleML contexts and
versions to annotated business-process tasks and a Regorous/PCL reasoner. Their
procedure accumulates effects, derives obligations in force, carries pending
obligations forward and checks fulfillment, violation and compensation. It
distinguishes all-trace compliance from existence of a compliant trace, and
conditions its reasoning claim on appropriate rules and complete annotations.
Omitting a trigger can prevent an obligation from appearing at all, producing a
deceptively clean result.

In the reported telecommunications pilot, roughly 100 paragraphs and associated
definitions became 176 normative statements; initial legal mapping took about
two weeks, followed by process-modeling workshops. These quantities show that
formalization and operational mapping are real work; they do not estimate the
time required by this bank. The proposal budgets compliance participation and
counts preparation in every comparison instead of treating gold rules as free
preprocessing.
% END SOURCE UNIT P-02-literature:06

% BEGIN REVISION ADDITION literature-addition
\section{A genuine citation can still support the wrong answer}
\label{sec:legal-ai-source-support}

The legal-AI evaluation literature adds a failure that a source-preservation
system alone cannot detect. An answer may link to an authentic document and
still misstate its consequence, use an inapplicable jurisdiction, or rely on a
position that has changed. The review therefore needs two questions: whether
the stated proposition is correct, and whether the cited authority supports it
for the activity and time under assessment.

\citet[sections 4--5 and appendix C]{legalrag2024} distinguish factual
correctness from legal groundedness when evaluating legal research tools. Their
2024 study used 202 queries, recorded product responses and manually applied
a published coding rubric. A source can be real yet fail to support the
proposition attributed to it. Their treatment of partially correct answers and
refusals also shows why a label must be read with its scoring definition. The
study's products, dates and United States legal questions do not estimate
LegalMath's error rate on Hong Kong circulars. Its useful contribution here is
the separation of answer correctness, citation support and incomplete answers.

\begin{figure}[H]
\centering
\TeachingCompare{Authentic document}{Supported proposition}{Applicable authority and date}
\caption{Three checks are needed before a citation can support a regulatory control.}
\label{fig:citation-three-checks}
\end{figure}

Consider a synthetic drafting error. A candidate states that an SFC request is
always required before a particular action, citing a paragraph that requires a
response \emph{when} the SFC requests it. The document exists and the quoted
words are accurate, but the candidate has changed the trigger. The audit should
retain the candidate's precise claim, the relevant sentence and its surrounding
conditions, and the reason the implication is unsupported. A quotation is useful
because another reviewer can inspect that comparison; quotation marks themselves
do not establish that the inference is sound.

\citet[sections 4.2--4.3 and appendix C]{legalfictions2024} distinguish
reference-based checks from a method that detects contradictions between model
responses. They inspect United States case-law tasks and validate the
contradiction labelling against human recoding of a sample. The distinction
matters to an ensemble: detecting a contradiction identifies an inconsistent
pair, while agreement leaves shared error possible. The workbench consequently
needs source support even when repeated answers are consistent. The paper's
published error estimates and its model-based labelling are not a substitute
for an independent Hong Kong reference process.

A citation audit must also record the edition actually read. A preprint, a
conference version, a later journal article and a changing software repository
may support different statements. If a later edition changes a theorem or a
benchmark denominator, retaining only the newest file destroys the explanation
of an earlier claim. The appropriate record binds the document bytes, section or
page, a short exact quotation, the manuscript proposition and the support
judgment. The full retained document allows examination beyond the short quote.

\section{Where the literature still leaves this product exposed}
\label{sec:remaining-literature-gaps}

The collection is strongest on executable legal languages, formal comparisons
and candidate-generation methods. Those contributions make the implementation
question concrete. They leave several questions on which the product's actual
banking value depends.

First, the product needs evidence about omissions across a \emph{complete}
regulatory packet, including incorporated definitions and attachments. A
benchmark whose input already supplies the decisive facts and formal rules
cannot answer whether a workbench discovers those facts and rules. The next
comparison should preserve independently inventoried source units and measure
which material units and alternatives are missed. Existing engineering cases
can remain as development checks without serving as the legal reference.

Second, comparison of language backends needs matched meanings. Before choosing
among RuleIR, Catala, a decision table or a normative language, give each the same
unknown facts, conflicting evidence, exceptions, dates and duties. Record where
an adapter cannot express the accepted meaning. A familiar syntax or a successful
example does not establish interchangeability. The restricted initial runtime
keeps this comparison manageable.

\begin{figure}[H]
\centering
\TeachingFlow{One reviewed banking question}{Same facts, authority and resource budget}{Compare consequential failures}
\caption{Coverage improves through matched questions, not by accumulating unrelated benchmark scores.}
\label{fig:literature-matched-question}
\end{figure}

Third, common-mode failure deserves a direct study. Changing the model vendor
while keeping the same defective extraction may leave the most serious omission
untouched. A useful experiment varies extraction, inventory and interpretation
separately and retains which activity first detected a defect. The numerical
dependence example in Chapter~\ref{ch:ensemble} explains the mechanism; it supplies
no measured correlation for this product.

Finally, the proposed human workflow needs independent evaluation. Reviewers
must be able to identify a plausible but unsupported source claim, notice a
minority reading and act on unresolved issues without an unmanageable queue.
The study in Chapter~\ref{ch:evaluation} therefore measures consequential errors
and review effort together. A legal-language theorem, model benchmark or
successful browser test does not answer that human question. These are defined
coverage gaps and research tasks, rather than a claim that no relevant work
exists beyond the retained literature.

% END REVISION ADDITION literature-addition

% BEGIN SOURCE UNIT M04:14
\section{A synthesis for the first implementation}
\label{merge:M04:14}

\label{sec:languages-synthesis}
The immediate product should keep a small executable core whose semantics are
fully specified. Around it, the interpretation record should be richer than
the executable language: it needs alternative readings, source versions,
definitions, unresolved qualitative judgments and supporting arguments.
Controlled templates give reviewers a readable view of formal structure.
Specialized analyses can then inspect declared fragments through adapters.

This choice avoids forcing every legal question into the first compiler.
A paragraph requiring judgment can become a reviewed assessment input or
manual procedure, with explicit evidence and responsibility. A temporal duty
can use an event profile whose limits are stated. A disputed definition can
remain in the interpretation set until review resolves it. The system can
still generate Java for a reviewed subset while keeping the broader
assessment visibly incomplete.

% BEGIN REVISION ADDITION teach-M04-14
\begin{figure}[H]
\centering
\TeachingFlow{One supported decision fragment}{One bounded event profile}{Measured additional capabilities}
\caption{The first implementation gives each borrowed mechanism a testable purpose.}\label{fig:teach-M04-14}
\end{figure}

% END REVISION ADDITION teach-M04-14
Every adapter needs a semantic contract. It declares supported constructs,
result categories, unknown and conflict behavior, vocabulary mapping,
resource bounds and how a returned witness is replayed. A paper's name or
a repository's existence does not supply that contract. A failed or
unsupported analysis is reported as such and cannot be counted as an
additional agreeing ensemble member.

The literature thus supplies complementary mechanisms: source-linked
authoring, precise exceptions, contextual alternatives, normative state,
symbolic consequence, path-directed testing and translation repair. The
new contribution proposed here is the controller that combines these
mechanisms without allowing one successful check to conceal another
unresolved question. Before specifying that controller, we need to explain
how competing interpretations are generated and defended. That is the task
of the next chapter.
% END SOURCE UNIT M04:14

